% Density revision 2026-10-04: answers below questions; compact related content.
% 4.x হাইপারলিংক in HTML
\section{HTML-এ হাইপারলিংক (Hyperlink)}

% Composition pass 2026-10-04
Home, Notice ও Result পেজ তৈরি আছে, কিন্তু এক পেজ থেকে অন্যটিতে যাওয়ার পথ দরকার।

\chFourMotivation{পেজে চলাচল $\rightarrow$ হাইপারলিংক $\rightarrow$ href-এ গন্তব্য ও ক্লিকযোগ্য লেখা/ছবি $\rightarrow$ কলেজের মেনু থেকে প্রয়োজনীয় পেজ খোলা।}


\begin{learningbox}
এই অংশ শেষে শিক্ষার্থী পারবে:
\begin{itemize}
\item হাইপারলিংক কী এবং ওয়েবপেজ চলাচলে এর প্রয়োজনীয়তা ব্যাখ্যা করতে।
\item \texttt{<a>} ট্যাগ ও \texttt{href} অ্যাট্রিবিউটের কাজ লিখতে।
\item অভ্যন্তরীণ, বাহ্যিক, বুকমার্ক, ইমেইল, ছবি, ফাইল ও নতুন ট্যাবের লিংক তৈরি করতে।
\item লিংকের লেখা, URL, \texttt{mailto:}, \texttt{target} ও \texttt{download} অ্যাট্রিবিউট ব্যবহার করতে।
\item পূর্ণ URL, আপেক্ষিক URL এবং নতুন ট্যাবে লিংক খোলার নিয়ম ব্যাখ্যা করতে।
\item HSC ব্যবহারিক অনুশীলন ও বোর্ড পরীক্ষার হাইপারলিংক-সম্পর্কিত প্রশ্নের উত্তর দিতে।
\end{itemize}
\end{learningbox}



HTML তালিকা শেখার পর হাইপারলিংক শেখা গুরুত্বপূর্ণ, কারণ একটি ওয়েবপেজের লেখা, তালিকা, ছবি বা বোতামের মাধ্যমে অন্য পেজ, ওয়েবসাইট, ফাইল বা একই পেজের নির্দিষ্ট অংশে যাওয়ার ব্যবস্থা হাইপারলিংক দিয়ে করা হয়। ওয়েবসাইটের চলাচল মেনু, ``Read more'', ``Download'', ``Contact us'', ছবি বোতাম - এসব হাইপারলিংকের বাস্তব উদাহরণ।

\subsection{হাইপারলিংক কী?}

\begin{definitionbox}{হাইপারলিংক (Hyperlink)}
ওয়েবপেজের যে লেখা, ছবি বা উপাদান নির্বাচন করলে অন্য পেজ, ওয়েবসাইট, ফাইল, ইমেইল ঠিকানা কিংবা একই পেজের নির্দিষ্ট অংশে যাওয়া যায়, তাকে হাইপারলিংক বলে।
\end{definitionbox}

হাইপারলিংক পেজ ও তথ্যের মধ্যে সংযোগ তৈরি করে। ব্যবহারকারী এই সংযোগ ধরে প্রয়োজনীয় তথ্যে পৌঁছান। তাই ওয়েবসাইটে সহজে চলাচলের জন্য অর্থবহ লিংক থাকা প্রয়োজন।



\subsection{অ্যাঙ্কর ট্যাগ: \texttt{<a>}}

\begin{definitionbox}{অ্যাঙ্কর ট্যাগ (Anchor Tag)}
HTML-এ হাইপারলিংক তৈরি করার জন্য \texttt{<a>} ট্যাগ ব্যবহার করা হয়। \texttt{<a>} ট্যাগকে অ্যাঙ্কর ট্যাগ বলা হয়।
\end{definitionbox}

অ্যাঙ্কর ট্যাগের সবচেয়ে গুরুত্বপূর্ণ অ্যাট্রিবিউট হলো \texttt{href}। \texttt{href} ছাড়া \texttt{<a>} ট্যাগ লিখলে সাধারণত লিংক কার্যকর হয় না, কারণ ব্রাউজার জানে না ক্লিক করলে কোথায় যেতে হবে।

\subsubsection{লেখার সাধারণ রূপ}

\begin{examplebox}{কোড}
\par\noindent\begin{minipage}{\linewidth}
\begin{verbatim}
<a href="url">Link Text</a>
\end{verbatim}
\end{minipage}\par
\end{examplebox}

\begin{center}
\begin{tabular}{|p{0.22\textwidth}|p{0.62\textwidth}|}
\hline
\textbf{অংশ} & \textbf{ব্যাখ্যা} \\
\hline
\texttt{<a>} & অ্যাঙ্কর ট্যাগ শুরু করে \\
\hline
\texttt{href} & লিংকের গন্তব্য বা ঠিকানা নির্ধারণ করে \\
\hline
\texttt{"url"} & যে পেজ, ওয়েবসাইট, ফাইল বা ঠিকানায় লিংক যাবে \\
\hline
\texttt{Link Text} & ব্যবহারকারী যে লেখায় ক্লিক করবে \\
\hline
\texttt{</a>} & অ্যাঙ্কর ট্যাগ শেষ করে \\
\hline
\end{tabular}
\end{center}

\begin{chFourWholeExample}{কোড ও ব্রাউজারে ফল}
\noindent
\begin{minipage}[t]{.50\linewidth}\vspace{0pt}

\begin{verbatim}
<a href="https://www.google.com">Google</a>
\end{verbatim}

\end{minipage}\hfill
\begin{minipage}[t]{.46\linewidth}\vspace{0pt}
\textbf{ফল:} ব্রাউজারে \textbf{Google} নামে একটি ক্লিকযোগ্য লিংক দেখা যাবে। এতে ক্লিক করলে Google ওয়েবসাইট খুলবে।
\end{minipage}
\chFourWhy{এখানে \texttt{href} অ্যাট্রিবিউট লিংকের গন্তব্য হিসেবে Google ওয়েবসাইটের ঠিকানা দিয়েছে এবং \texttt{Google} হলো ক্লিকযোগ্য লিংকের লেখা।}
\end{chFourWholeExample}

\subsection{\texttt{href} অ্যাট্রিবিউট}

\begin{definitionbox}{\texttt{href}}
\texttt{href} অ্যাট্রিবিউট লিংকের গন্তব্য ঠিকানা নির্ধারণ করে। Hyperlink Reference কথাটি দিয়ে এর কাজ বোঝানো হয়; কোডে অ্যাট্রিবিউটের নাম \texttt{href}।
\end{definitionbox}

\begin{keypointbox}{মনে রাখো}
\texttt{href} ছাড়া \texttt{<a>} এলিমেন্ট হাইপারলিংক তৈরি করে না; এটি লিংকের জন্য স্থানধারক হিসেবে থাকতে পারে।
\end{keypointbox}


\begin{keypointbox}{লিংকের লেখা বাছাইয়ের নিয়ম}
\texttt{Click here} এর চেয়ে \texttt{Download Routine}, \texttt{Visit Google} বা \texttt{Go to Contact} - এ ধরনের অর্থবহ লিংকের লেখা ব্যবহার করা ভালো।
\end{keypointbox}

\subsection{হাইপারলিংকের প্রকারভেদ}

গন্তব্য, ক্লিকযোগ্য অংশ ও খোলার পদ্ধতি অনুযায়ী লিংকের নিচের রূপগুলো পরিচিত। এগুলো পরস্পর বিচ্ছিন্ন শ্রেণি নয়; একই বাহ্যিক লিংক ছবিতে দেওয়া যেতে পারে এবং নতুন ট্যাবেও খুলতে পারে:

\begin{enumerate}
\item বাহ্যিক লিংক
\item অভ্যন্তরীণ লিংক
\item একই পেজ বা বুকমার্ক লিংক
\item ইমেইল লিংক
\item ছবিতে লিংক
\item ডাউনলোড বা ফাইলের লিংক
\item নতুন ট্যাবের লিংক
\end{enumerate}

\noindent
\begin{minipage}[t]{.48\linewidth}\vspace{0pt}
\begin{definitionbox}{বাহ্যিক লিংক (External Link)}
বাহ্যিক লিংক হলো এমন হাইপারলিংক, যা বর্তমান ওয়েবসাইট বা বর্তমান পেজ থেকে অন্য কোনো ওয়েবসাইটে নিয়ে যায়।
\end{definitionbox}
\end{minipage}\hfill
\begin{minipage}[t]{.48\linewidth}\vspace{0pt}
\begin{definitionbox}{অভ্যন্তরীণ লিংক (Internal Link)}
অভ্যন্তরীণ লিংক হলো এমন হাইপারলিংক, যা একই ওয়েবসাইটের একটি পেজ থেকে অন্য পেজে নিয়ে যায়।
\end{definitionbox}
\end{minipage}\par

\subsection{বাহ্যিক লিংক}



বাহ্যিক লিংক লিখতে সাধারণত ওয়েবসাইট ঠিকানার আগে \texttt{http://} বা \texttt{https://} লেখা হয়।

\subsubsection{উদাহরণ}

% Split revision: code and its explanation share one row.
\par\noindent
\begin{minipage}[t]{.54\linewidth}\vspace{0pt}\small
\textbf{\color{chFourTeal}কোড}\par
\begin{verbatim}
<a href="https://www.google.com">Visit Google</a>
\end{verbatim}

\par\smallskip
\end{minipage}\hfill
\begin{minipage}[t]{.42\linewidth}\vspace{0pt}\small
\textbf{\color{chFourTeal}ফল ও ব্যাখ্যা}\par
\textbf{ফল:} ব্রাউজারে \textbf{Visit Google} লেখা দেখা যাবে। এতে ক্লিক করলে Google ওয়েবসাইট খুলবে।

\textbf{ব্যাখ্যা:} এখানে লিংক বর্তমান ওয়েবপেজের বাইরে অন্য ওয়েবসাইটে যাচ্ছে, তাই এটি বাহ্যিক লিংক।
\par\smallskip
\end{minipage}\par




\subsection{অভ্যন্তরীণ লিংক}



একটি কলেজের ওয়েবসাইটে \texttt{Home} পেজ থেকে \texttt{About}, \texttt{Notice}, \texttt{Result} বা \texttt{Contact} পেজে যাওয়ার জন্য অভ্যন্তরীণ লিংক ব্যবহার করা হয়।

\subsubsection{উদাহরণ}

\begin{chFourWholeExample}{কোড ও ব্রাউজারে ফল}
\noindent
\begin{minipage}[t]{.50\linewidth}\vspace{0pt}

\begin{verbatim}
<a href="about.html">About Us</a>
<a href="notice.html">Notice</a>
<a href="contact.html">Contact</a>
\end{verbatim}

\end{minipage}\hfill
\begin{minipage}[t]{.46\linewidth}\vspace{0pt}
\textbf{ফল:} ব্রাউজারে \textbf{About Us}, \textbf{Notice}, \textbf{Contact} নামে তিনটি লিংক দেখা যাবে।
\end{minipage}
\chFourWhy{এগুলো একই ওয়েবসাইটের স্থানীয় পেজে লিংক করে। তাই এগুলো অভ্যন্তরীণ লিংক।}
\end{chFourWholeExample}

\subsubsection{পূর্ণাঙ্গ উদাহরণ: \texttt{index.html}}

% Split revision: code and its explanation share one row.
\par\noindent
\begin{minipage}[t]{.50\linewidth}\vspace{0pt}\small
\textbf{\color{chFourTeal}কোড}\par
\begin{verbatim}
<!DOCTYPE html>
<html>
<body>
  <h1>ABC College</h1>

  <a href="about.html">About</a>
  <a href="notice.html">Notice</a>
  <a href="contact.html">Contact</a>
</body>
</html>
\end{verbatim}

\par\smallskip
\end{minipage}\hfill
\begin{minipage}[t]{.46\linewidth}\vspace{0pt}\small

\end{minipage}\par


% Original applied explanation; see Chapter04/SOURCES.md.
\subsubsection{Notice লিংক ভুল ফোল্ডারে যায় কেন?}
ধরা যাক পেজটি \texttt{/students/index.html}-এ আছে। সেখানে \texttt{href="notice.html"} দিলে সাধারণ অবস্থায় \texttt{/students/notice.html} খোঁজা হবে; \texttt{../notice.html} দিলে এক স্তর ওপরের \texttt{/notice.html} খোঁজা হবে। তাই আপেক্ষিক পথ ঠিক কি না, তা বর্তমান পেজের ভিত্তি ঠিকানা ধরে বিচার করবে।

দর্শককে গন্তব্য বোঝাতে \textbf{ভর্তির বিজ্ঞপ্তি পড়ুন} ধরনের লিংক-লেখা ব্যবহার করবে। একই সাইটে সব লিংকে \textbf{এখানে ক্লিক করুন} লিখলে কোন লিংক কোন তথ্য দেবে তা বোঝা কঠিন হয়।


\subsection{পূর্ণ URL ও আপেক্ষিক URL}

\begin{definitionbox}{পূর্ণ ঠিকানা (Absolute URL)}
পূর্ণ URL হলো সম্পূর্ণ ওয়েব ঠিকানা, যেখানে প্রোটোকল, ডোমেইন এবং প্রয়োজনীয় পথ একসঙ্গে লেখা থাকে।
\end{definitionbox}

\begin{definitionbox}{আপেক্ষিক ঠিকানা (Relative URL)}
যে ঠিকানা একটি ভিত্তি ঠিকানার সঙ্গে মিলিয়ে সম্পূর্ণ করা হয়, তাকে \textbf{আপেক্ষিক URL} বলে। একই ওয়েবসাইটের পেজ বা ফাইলের ক্ষেত্রে প্রায়ই শুধু আপেক্ষিক পথ লেখা হয়; যেমন \texttt{about.html}।
\end{definitionbox}

\begin{center}
\begin{tabular}{|p{0.28\textwidth}|p{0.30\textwidth}|p{0.28\textwidth}|}
\hline
\textbf{বিষয়} & \textbf{পূর্ণ URL} & \textbf{আপেক্ষিক URL} \\
\hline
ঠিকানা & সম্পূর্ণ ঠিকানা & ভিত্তি ঠিকানার সঙ্গে মিলিয়ে সম্পূর্ণ করতে হয় \\
\hline
প্রোটোকল & থাকে & সাধারণত থাকে না \\
\hline
ব্যবহার & একই বা অন্য সাইটের পূর্ণ ঠিকানা & ভিত্তি ঠিকানা অনুযায়ী গন্তব্য \\
\hline
উদাহরণ & \texttt{https://google.com} & \texttt{about.html} \\
\hline
\end{tabular}
\end{center}

\subsubsection{লিংক থেকে গন্তব্য বোঝার ধাপ}
একটি আপেক্ষিক ঠিকানা বুঝতে বর্তমান পেজের ভিত্তি ঠিকানাও লাগে। নিচের উদাহরণে ভিত্তি ঠিকানা \texttt{https://college.example/} ধরা হয়েছে।
\begin{center}
\resizebox{.96\linewidth}{!}{%
\begin{tikzpicture}[node distance=7mm,
  linkstep/.style={draw=chFourTeal,rounded corners,fill=chFourTint,
    text width=38mm,minimum height=18mm,align=center,font=\small},
  flow/.style={-{Stealth[length=2mm]},thick,draw=chFourTeal}]
\node[linkstep] (a) {ক্লিকযোগ্য লেখা\\\texttt{About Us}};
\node[linkstep,right=of a] (b) {লিংকের ঠিকানা\\\texttt{href="about.html"}};
\node[linkstep,right=of b] (c) {ভিত্তির সঙ্গে মিলিয়ে\\গন্তব্য পেজ খোলা};
\draw[flow] (a)--(b); \draw[flow] (b)--(c);
\end{tikzpicture}}
\end{center}
ফলে গন্তব্য হয় \texttt{https://college.example/about.html}। ক্লিকযোগ্য লেখা বদলালেও \texttt{href} একই থাকলে গন্তব্য বদলায় না; \texttt{href} বদলালে গন্তব্য বদলায়।



\subsection{একই পেজের লিংক বা বুকমার্ক লিংক}

\begin{definitionbox}{একই পেজের লিংক (Same Page / Bookmark Link)}
একই পেজ লিংক হলো এমন হাইপারলিংক, যার মাধ্যমে একই ওয়েবপেজের নির্দিষ্ট অংশে যাওয়া যায়।
\end{definitionbox}

বড় ওয়েবপেজ, নিবন্ধ, সূচিপত্র বা বহুল জিজ্ঞাসিত প্রশ্নের অংশে একই পেজের লিংক ব্যবহার করলে ব্যবহারকারী দ্রুত প্রয়োজনীয় অংশে যেতে পারে।

\subsubsection{উদাহরণ}

% Split revision: code and its explanation share one row.
\par\noindent
\begin{minipage}[t]{.54\linewidth}\vspace{0pt}\small
\textbf{\color{chFourTeal}কোড}\par
\begin{verbatim}
<a href="#html">Go to HTML Section</a>

<h2 id="html">HTML Section</h2>
<p>HTML is used to create web page structure.</p>
\end{verbatim}

\par\smallskip
\end{minipage}\hfill
\begin{minipage}[t]{.42\linewidth}\vspace{0pt}\small
\textbf{\color{chFourTeal}ফল ও ব্যাখ্যা}\par
\textbf{ফল:} ব্রাউজারে \textbf{Go to HTML Section} নামে লিংক দেখা যাবে। এতে ক্লিক করলে একই পেজের \textbf{HTML Section} অংশে চলে যাবে।

\textbf{ব্যাখ্যা:} \texttt{href="\#html"} মানে পেজের \texttt{id="html"} অংশে যাওয়া। এখানে \texttt{\#} দিয়ে গন্তব্য id নির্দেশ করা হয়েছে।
\par\smallskip
\end{minipage}\par




\subsection{ইমেইল লিংক}

\begin{definitionbox}{ইমেইল লিংক (Email Link)}
ইমেইল লিংক হলো এমন হাইপারলিংক, যার মাধ্যমে নির্দিষ্ট ঠিকানায় বার্তা লেখার জন্য ইমেইল অ্যাপ্লিকেশন খোলার অনুরোধ করা হয়। এটি কাজ করতে উপযুক্ত অ্যাপ বা ওয়েবমেইল ব্যবস্থা থাকতে হয়; ক্লিক করলেই বার্তা পাঠানো হয় না।
\end{definitionbox}

ইমেইল লিংক তৈরি করতে \texttt{href} অ্যাট্রিবিউটের মান হিসেবে \texttt{mailto:} ব্যবহার করা হয়।

\subsubsection{সিনট্যাক্স}

\begin{examplebox}{কোড}
\par\noindent\begin{minipage}{\linewidth}
\begin{verbatim}
<a href="mailto:info@example.com">Send Email</a>
\end{verbatim}
\end{minipage}\par
\end{examplebox}

\subsubsection{উদাহরণ}

% Split revision: code and its explanation share one row.
\par\noindent
\begin{minipage}[t]{.54\linewidth}\vspace{0pt}\small
\textbf{\color{chFourTeal}কোড}\par
\begin{verbatim}
<a href="mailto:principal@abccollege.edu.bd">Email Principal</a>
\end{verbatim}

\par\smallskip
\end{minipage}\hfill
\begin{minipage}[t]{.42\linewidth}\vspace{0pt}\small
\textbf{\color{chFourTeal}ফল ও ব্যাখ্যা}\par
\textbf{ফল:} ব্রাউজারে \textbf{Email Principal} লিংক দেখা যাবে। এতে ক্লিক করলে ইমেইল অ্যাপ্লিকেশন খুলতে পারে এবং প্রাপক হিসেবে \texttt{principal@abccollege.edu.bd} ঠিকানা বসবে।

\textbf{ব্যাখ্যা:} \texttt{mailto:} ব্রাউজারকে জানায় যে লিংকটি ইমেইল ঠিকানার জন্য।
\par\smallskip
\end{minipage}\par




\subsection{ছবিতে লিংক}

\begin{definitionbox}{ছবিতে লিংক (Image Link)}
ছবিতে লিংক হলো এমন হাইপারলিংক যেখানে লিংকের লেখার পরিবর্তে ছবি ব্যবহার করা হয়। ছবিতে ক্লিক করলে নির্দিষ্ট পেজ বা ওয়েবসাইট খোলে।
\end{definitionbox}

ছবিতে লিংক তৈরি করতে \texttt{<a>} ট্যাগের ভিতরে \texttt{<img>} ট্যাগ বসানো হয়।

\subsubsection{সিনট্যাক্স}

\begin{examplebox}{কোড}
\par\noindent\begin{minipage}{\linewidth}
\begin{verbatim}
<a href="url">
  <img src="image.jpg" alt="Image">
</a>
\end{verbatim}
\end{minipage}\par
\end{examplebox}

\subsubsection{উদাহরণ}

% Split revision: code and its explanation share one row.
\par\noindent
\begin{minipage}[t]{.54\linewidth}\vspace{0pt}\small
\textbf{\color{chFourTeal}কোড}\par
\begin{verbatim}
<a href="https://www.ictacademy.com">
  <img src="ict.jpg" alt="ICT Academy" width="200">
</a>
\end{verbatim}

\par\smallskip
\end{minipage}\hfill
\begin{minipage}[t]{.42\linewidth}\vspace{0pt}\small
\textbf{\color{chFourTeal}ফল ও ব্যাখ্যা}\par
\textbf{ফল:} ব্রাউজারে \texttt{ict.jpg} ছবি দেখা যাবে। ছবিতে ক্লিক করলে নির্দিষ্ট ওয়েবসাইট খুলবে।

\textbf{ব্যাখ্যা:} এখানে ক্লিকযোগ্য অংশটি লেখা নয়, ছবি। তাই \texttt{<a>} ট্যাগের ভিতরে \texttt{<img>} ট্যাগ ব্যবহার করা হয়েছে।
\par\smallskip
\end{minipage}\par




\subsection{ডাউনলোড বা ফাইলের লিংক}

\begin{definitionbox}{ফাইলের লিংক (File Link)}
যে হাইপারলিংক কোনো PDF, ছবি, ডকুমেন্ট বা অন্য ফাইল খুলে/ডাউনলোড করার জন্য ব্যবহৃত হয়, তাকে ফাইলের লিংক বলা যায়।
\end{definitionbox}

\subsubsection{উদাহরণ}

\begin{chFourWholeExample}{কোড ও ব্রাউজারে ফল}
\noindent
\begin{minipage}[t]{.50\linewidth}\vspace{0pt}

\begin{verbatim}
<a href="routine.pdf">Download Routine</a>
<a href="syllabus.pdf">Download Syllabus</a>
\end{verbatim}

\end{minipage}\hfill
\begin{minipage}[t]{.46\linewidth}\vspace{0pt}
\textbf{ফল:} ব্রাউজারে \textbf{Download Routine} এবং \textbf{Download Syllabus} নামে লিংক দেখা যাবে।
\end{minipage}
\chFourWhy{লিংকে ক্লিক করলে ব্রাউজার ফাইল খুলতে পারে অথবা ডাউনলোড করতে পারে। এটি ব্রাউজারের সেটিং ও ফাইলের ধরনের উপর নির্ভর করে।}
\end{chFourWholeExample}

\subsubsection{\texttt{download} অ্যাট্রিবিউট}

% Split revision: code and its explanation share one row.
\par\noindent
\begin{minipage}[t]{.54\linewidth}\vspace{0pt}\small
\textbf{\color{chFourTeal}কোড}\par
\begin{verbatim}
<a href="routine.pdf" download>Download Routine</a>
\end{verbatim}

\par\smallskip
\end{minipage}\hfill
\begin{minipage}[t]{.42\linewidth}\vspace{0pt}\small
\textbf{\color{chFourTeal}ফল ও ব্যাখ্যা}\par
\textbf{ফল:} \textbf{Download Routine} লিংকে ক্লিক করলে ফাইল ডাউনলোড হওয়ার নির্দেশ পায়।

\textbf{ব্যাখ্যা:} \texttt{download} অ্যাট্রিবিউট ব্রাউজারকে ফাইল ডাউনলোড করতে উৎসাহিত করে। তবে কিছু ব্রাউজার বা সার্ভারের নীতির কারণে আচরণ ভিন্ন হতে পারে।
\par\smallskip
\end{minipage}\par


\subsection{নতুন ট্যাবের লিংক}

\begin{definitionbox}{নতুন ট্যাবের লিংক (New Tab Link)}
\texttt{target="\_blank"} অ্যাট্রিবিউট ব্যবহার করলে লিংক নতুন ট্যাব বা নতুন উইন্ডোতে খোলে।
\end{definitionbox}

\subsubsection{উদাহরণ}

% Split revision: code and its explanation share one row.
\par\noindent
\begin{minipage}[t]{.54\linewidth}\vspace{0pt}\small
\textbf{\color{chFourTeal}কোড}\par
\begin{verbatim}
<a href="https://www.google.com" target="_blank">Open Google</a>
\end{verbatim}

\par\smallskip
\end{minipage}\hfill
\begin{minipage}[t]{.42\linewidth}\vspace{0pt}\small
\textbf{\color{chFourTeal}ফল ও ব্যাখ্যা}\par
\textbf{ফল:} ব্রাউজারে \textbf{Open Google} লিংক দেখা যাবে। এতে ক্লিক করলে Google নতুন ট্যাবে খুলবে।

\textbf{ব্যাখ্যা:} \texttt{target="\_blank"} গন্তব্য পেজ নতুন ট্যাব/উইন্ডোতে খোলার নির্দেশ দেয়।
\par\smallskip
\end{minipage}\par


\subsubsection{পরিচিত \texttt{target} মান}

\begin{center}
\begin{tabular}{|p{0.20\textwidth}|p{0.60\textwidth}|}
\hline
\textbf{মান} & \textbf{কাজ} \\
\hline
\texttt{\_self} & একই ট্যাব/উইন্ডোতে খোলে; স্বাভাবিক আচরণ \\
\hline
\texttt{\_blank} & নতুন ট্যাব/উইন্ডোতে খোলে \\
\hline
\texttt{\_parent} & বাইরের ফ্রেমে খোলে \\
\hline
\texttt{\_top} & পুরো ব্রাউজার উইন্ডোতে খোলে \\
\hline
\end{tabular}
\end{center}

\begin{keypointbox}{আধুনিক ব্যবহারের নিয়ম}
\texttt{rel="noopener"} নতুন পেজকে মূল পেজের \texttt{window.opener} সংযোগ পাওয়া থেকে আটকায়। \texttt{noreferrer} গন্তব্যে উৎস পেজের Referer তথ্য পাঠানো বন্ধ করে। বর্তমান HTML-এ \texttt{target="\_blank"} সাধারণত নিজেই opener সংযোগ বিচ্ছিন্ন করে; নিচের উদাহরণে দুটি উদ্দেশ্য স্পষ্ট করে লেখা হয়েছে।
\end{keypointbox}

\begin{examplebox}{কোড}
\par\noindent\begin{minipage}{\linewidth}
\begin{verbatim}
<a href="https://www.google.com" target="_blank" rel="noopener noreferrer">
  Open Google
</a>
\end{verbatim}
\end{minipage}\par
\end{examplebox}

\subsection{লিংকের স্বাভাবিক রং}

সাধারণভাবে ব্রাউজার হাইপারলিংকের অবস্থা অনুযায়ী ভিন্ন রং দেখায়। নিজস্ব CSS না থাকলে আগে দেখা হয়নি এমন লিংক সাধারণত \textbf{নীল} হয়।

\begin{center}
\begin{tabular}{|p{0.35\textwidth}|p{0.35\textwidth}|}
\hline
\textbf{লিংকের অবস্থা} & \textbf{সাধারণ রং} \\
\hline
আগে দেখা হয়নি এমন লিংক & নীল \\
\hline
আগে দেখা হয়েছে এমন লিংক & বেগুনি \\
\hline
ক্লিক করার সময়কার লিংক & লাল \\
\hline
\end{tabular}
\end{center}

\begin{boardbox}{পরীক্ষায় মনে রাখবে}
স্বাভাবিক হাইপারলিংকের রং সাধারণত \textbf{নীল}।
\end{boardbox}

\subsection{হাইপারলিংকের সুবিধা}

হাইপারলিংক ওয়েবপেজকে সম্পর্কযুক্ত করে এবং ব্যবহারকারীর চলাচলের সুযোগ বাড়ায়। এর প্রধান সুবিধাগুলো হলো:

\par\noindent
\begin{minipage}[t]{.48\linewidth}\vspace{0pt}
\begin{enumerate}
\item এক পেজ থেকে অন্য পেজে দ্রুত যাওয়া যায়।
\item সম্পর্কিত তথ্য সহজে দেখা যায়।
\item ওয়েবসাইটের পেজগুলোর মধ্যে চলাচল সহজ হয়।
\item বাহ্যিক সম্পদ যুক্ত করা যায়।
\end{enumerate}
\end{minipage}\hfill
\begin{minipage}[t]{.48\linewidth}\vspace{0pt}
\begin{enumerate}
\item PDF, ছবি বা ডকুমেন্ট খুলে/ডাউনলোড করা যায়।
\item ইমেইল যোগাযোগ সহজ হয়।
\item একই পেজের নির্দিষ্ট অংশে দ্রুত যাওয়া যায়।
\end{enumerate}
\end{minipage}\par


\subsection{হাইপারলিংকের লেখার রূপ}

\begingroup
\setlength{\LTpre}{2pt}\setlength{\LTpost}{2pt}
\begin{longtable}{|p{0.23\textwidth}|p{0.57\textwidth}|}
\hline
\textbf{লিংকের ধরন} & \textbf{সিনট্যাক্স} \\
\hline
\endfirsthead
\hline
\textbf{লিংকের ধরন} & \textbf{সিনট্যাক্স} \\
\hline
\endhead
বাহ্যিক লিংক & {\footnotesize\ttfamily <a href="https://example.com">}\\
& {\footnotesize\ttfamily Visit</a>} \\
\hline
অভ্যন্তরীণ লিংক & {\footnotesize\ttfamily <a href="about.html">About</a>} \\
\hline
একই পেজ লিংক & {\footnotesize\ttfamily <a href="\#section1">Go</a>} \\
\hline
ইমেইল লিংক & {\footnotesize\ttfamily <a href="mailto:info@example.com">}\\
& {\footnotesize\ttfamily Email</a>} \\
\hline
ছবিতে লিংক & {\footnotesize\ttfamily <a href="url">}\\
& {\footnotesize\ttfamily <img src="photo.jpg"></a>} \\
\hline
ফাইল লিংক & {\footnotesize\ttfamily <a href="routine.pdf">}\\
& {\footnotesize\ttfamily Download</a>} \\
\hline
নতুন ট্যাব লিংক & {\footnotesize\ttfamily <a href="https://example.com"}\\
& {\footnotesize\ttfamily target="\_blank">Open</a>} \\
\hline
\end{longtable}
\endgroup


\subsection{পূর্ণাঙ্গ উদাহরণ}

\Needspace{6\baselineskip}
\par\noindent{\color{chFourTeal}\bfseries\small কোড: বাম কলাম শেষে ডান কলাম পড়বে}\par
\begin{multicols}{2}
\begin{verbatim}
<h1>ABC College Website</h1>

<h2>Internal Links</h2>
<a href="about.html">About</a> |
<a href="notice.html">Notice</a> |
<a href="contact.html">Contact</a>

<h2>External Link</h2>
<a href="https://www.google.com"
   target="_blank">
  Visit Google
</a>

<h2>Email Link</h2>
<a href="mailto:info@abccollege.edu.bd">
  Email Us
</a>

<h2>Image Link</h2>
<a href="https://www.ictacademy.com">
  <img src="ict.jpg" alt="ICT Academy" width="200">
</a>

<h2>Same Page Link</h2>
<a href="#result">Go to Result Section</a>

<h2 id="result">Result Section</h2>
<p>HSC result information will be published here.</p>
\end{verbatim}
\end{multicols}


\begin{examplebox}{ফল কী হবে?}
ব্রাউজারে দেখা যাবে:
\par\noindent
\begin{minipage}[t]{.48\linewidth}\vspace{0pt}
\begin{itemize}
\item \textbf{ABC College Website} শিরোনাম।
\item \texttt{About}, \texttt{Notice}, \texttt{Contact} অভ্যন্তরীণ লিংক।
\item \texttt{Google} বাহ্যিক লিংক।
\end{itemize}
\end{minipage}\hfill
\begin{minipage}[t]{.48\linewidth}\vspace{0pt}
\begin{itemize}
\item \texttt{Email Us} ইমেইল লিংক।
\item একটি ছবি; ছবিতে ক্লিক করলে নির্দিষ্ট ওয়েবসাইট খুলবে।
\item \texttt{Go to Result Section} লিংক; এতে ক্লিক করলে একই পেজের \texttt{Result Section} অংশে যাবে।
\end{itemize}
\end{minipage}\par

\end{examplebox}

\subsection{সাধারণ ভুল}

\begin{mistakebox}{সাধারণ ভুল}
\par\noindent
\begin{minipage}[t]{.48\linewidth}\vspace{0pt}
\begin{itemize}
\item \texttt{href} অ্যাট্রিবিউট না দেওয়া।
\item \texttt{href}-এর জায়গায় \texttt{src} লেখা।
\item বাহ্যিক লিংকে \texttt{https://} না দেওয়া।
\item শেষের \texttt{</a>} ট্যাগ না দেওয়া।
\item \texttt{<a>} ট্যাগের ভিতরে ক্লিক করার লেখা বা ছবি না রাখা।
\end{itemize}
\end{minipage}\hfill
\begin{minipage}[t]{.48\linewidth}\vspace{0pt}
\begin{itemize}
\item ছবিতে লিংক করতে \texttt{<img>} ট্যাগ \texttt{<a>} ট্যাগের বাইরে রাখা।
\item ইমেইল লিংকে \texttt{mailto:} না লেখা।
\item একই পেজের লিংকে \texttt{\#} না দেওয়া।
\item \texttt{id} এবং \texttt{href="\#id"} মান মিল না রাখা।
\end{itemize}
\end{minipage}\par

\end{mistakebox}

\subsubsection{ভুল}

% Split revision: code and its explanation share one row.
\par\noindent
\begin{minipage}[t]{.50\linewidth}\vspace{0pt}\small
\textbf{\color{chFourTeal}কোড}\par
\begin{verbatim}
<a>Google</a>
\end{verbatim}

\par\smallskip
\end{minipage}\hfill
\begin{minipage}[t]{.46\linewidth}\vspace{0pt}\small
\textbf{\color{chFourTeal}ব্যাখ্যা}\par
এখানে \texttt{href} নেই, তাই এটি কার্যকর হাইপারলিংক হবে না।
\end{minipage}\par
\subsubsection{সঠিক}

\begin{examplebox}{কোড}
\par\noindent\begin{minipage}{\linewidth}
\begin{verbatim}
<a href="https://www.google.com">Google</a>
\end{verbatim}
\end{minipage}\par
\end{examplebox}

\subsection{সংক্ষিপ্ত অনুশীলন}

% Each question keeps its answer beneath it in the same column.
\begin{multicols}{2}
\noindent\begin{minipage}[t]{\linewidth}\vspace{0pt}
\chFourReviewQuestion{\textbf{হাইপারলিংক কী?}}
\end{minipage}\par\smallskip
\noindent\begin{minipage}[t]{\linewidth}\vspace{0pt}
\chFourReviewQuestion{\textbf{হাইপারলিংক তৈরি করতে কোন ট্যাগ ব্যবহৃত হয়?}}
\end{minipage}\par\smallskip
\noindent\begin{minipage}[t]{\linewidth}\vspace{0pt}
\chFourReviewQuestion{\textbf{\texttt{<a>} ট্যাগের প্রয়োজনীয় অ্যাট্রিবিউট কোনটি?}}
\end{minipage}\par\smallskip
\noindent\begin{minipage}[t]{\linewidth}\vspace{0pt}
\chFourReviewQuestion{\textbf{\texttt{href} এর পূর্ণ অর্থ কী?}}
\end{minipage}\par\smallskip
\noindent\begin{minipage}[t]{\linewidth}\vspace{0pt}
\chFourReviewQuestion{\textbf{হাইপারলিংকের সিনট্যাক্স কী?}}
\end{minipage}\par\smallskip
\noindent\begin{minipage}[t]{\linewidth}\vspace{0pt}
\chFourReviewQuestion{\textbf{ইমেইল লিংক তৈরি করতে কী ব্যবহার করা হয়?}}
\end{minipage}\par\smallskip
\noindent\begin{minipage}[t]{\linewidth}\vspace{0pt}
\chFourReviewQuestion{\textbf{ছবিতে লিংক কীভাবে তৈরি করা হয়?}}
\end{minipage}\par\smallskip
\noindent\begin{minipage}[t]{\linewidth}\vspace{0pt}
\chFourReviewQuestion{\textbf{নতুন ট্যাবে লিংক খুলতে কোন অ্যাট্রিবিউট ব্যবহার হয়?}}
\end{minipage}\par\smallskip
\end{multicols}

\subsection{বোর্ডভিত্তিক সৃজনশীল প্রশ্ন}

\begin{boardbox}{সৃজনশীল প্রশ্ন ১}
একটি কলেজ ওয়েবসাইটে \texttt{Home}, \texttt{Notice}, \texttt{Result}, \texttt{Email Us} এবং \texttt{Open Google} লিংক আছে। \texttt{Notice} ও \texttt{Result} একই ওয়েবসাইটের পেজে নিয়ে যায়, \texttt{Email Us} ইমেইল অ্যাপ্লিকেশন খোলার অনুরোধ করে, আর \texttt{Open Google} নতুন ট্যাবে যায়।
\end{boardbox}



\begin{boardbox}{সৃজনশীল প্রশ্ন ২}
একটি ওয়েবপেজে \texttt{Go to Admission Section} লিংক আছে, যা একই পেজের ভর্তি অংশে নিয়ে যায়। আরেকটি কোর্স পরিচিতিপত্র ছবিতে ক্লিক করলে সংশ্লিষ্ট কোর্স পেজ খোলে।
\end{boardbox}



\Needspace{10\baselineskip}
\subsection{পরীক্ষাভিত্তিক বহুনির্বাচনি প্রশ্ন}

\begin{multicols}{2}
\begin{enumerate}[leftmargin=*,itemsep=5pt]
\item \begin{minipage}[t]{\linewidth}\vspace{0pt}
HTML-এ হাইপারলিংক তৈরি করতে কোন ট্যাগ ব্যবহৃত হয়?
\par ক) \texttt{<p>}
\par খ) \texttt{<a>}
\par গ) \texttt{<li>}
\par ঘ) \texttt{<img>}
\end{minipage}

\item \begin{minipage}[t]{\linewidth}\vspace{0pt}
\texttt{<a>} ট্যাগের প্রয়োজনীয় অ্যাট্রিবিউট কোনটি?
\par ক) \texttt{src}
\par খ) \texttt{href}
\par গ) \texttt{alt}
\par ঘ) \texttt{border}
\end{minipage}

\item \begin{minipage}[t]{\linewidth}\vspace{0pt}
\texttt{href} এর অর্থ কী?
\par ক) Hot Reference
\par খ) Home Reference
\par গ) Hyperlink Reference
\par ঘ) Hyperlink Replace
\end{minipage}

\item \begin{minipage}[t]{\linewidth}\vspace{0pt}
হাইপারলিংকের সঠিক সিনট্যাক্স কোনটি?
\par ক) \texttt{<a url="link">Text</a>}
\par খ) \texttt{<a href="url">Link Text</a>}
\par গ) \texttt{<href="url">Text</href>}
\par ঘ) \texttt{<link href="url">Text</link>}
\end{minipage}

\item \begin{minipage}[t]{\linewidth}\vspace{0pt}
ইমেইল ঠিকানার সঙ্গে লিংক তৈরি করতে কোনটি ব্যবহার করা হয়?
\par ক) \texttt{email:}
\par খ) \texttt{mailto:}
\par গ) \texttt{src:}
\par ঘ) \texttt{target:}
\end{minipage}

\item \begin{minipage}[t]{\linewidth}\vspace{0pt}
ছবিতে লিংক তৈরি করতে কী করতে হয়?
\par ক) \texttt{<img>} ট্যাগের ভিতরে \texttt{<p>} ট্যাগ বসাতে হয়
\par খ) \texttt{<a>} ট্যাগের ভিতরে \texttt{<img>} ট্যাগ বসাতে হয়
\par গ) \texttt{<table>} ট্যাগের ভিতরে \texttt{<a>} ট্যাগ বসাতে হয়
\par ঘ) \texttt{<body>} ট্যাগ বাদ দিতে হয়
\end{minipage}

\item \begin{minipage}[t]{\linewidth}\vspace{0pt}
স্বাভাবিক আগে দেখা হয়নি এমন হাইপারলিংকের রং সাধারণত কী?
\par ক) লাল
\par খ) সবুজ
\par গ) নীল
\par ঘ) গোলাপি
\end{minipage}

\item \begin{minipage}[t]{\linewidth}\vspace{0pt}
\texttt{target="\_blank"} অ্যাট্রিবিউটের কাজ কী?
\par ক) লিংক মুছে ফেলা
\par খ) নতুন ট্যাব/উইন্ডোতে খোলা
\par গ) ফাইল ডাউনলোড করা
\par ঘ) একই পেজ নির্দিষ্ট অংশে যাওয়া
\end{minipage}

\item \begin{minipage}[t]{\linewidth}\vspace{0pt}
\texttt{download} অ্যাট্রিবিউটের কাজ কী?
\par ক) লিংক লেখার রং পরিবর্তন করা
\par খ) ফাইল ডাউনলোডের অনুরোধ করা
\par গ) ছবিতে লিংক তৈরি করা
\par ঘ) ইমেইল খোলা
\end{minipage}

\item \begin{minipage}[t]{\linewidth}\vspace{0pt}
\texttt{href="\#section"} দ্বারা কোন ধরনের লিংক তৈরি হয়?
\par ক) বাহ্যিক লিংক
\par খ) অভ্যন্তরীণ লিংক
\par গ) একই পেজ/বুকমার্ক লিংক
\par ঘ) ফাইল লিংক
\end{minipage}

\item \begin{minipage}[t]{\linewidth}\vspace{0pt}
\texttt{<a>} ট্যাগে \texttt{href} না থাকলে সাধারণত কী হয়?
\par ক) লিংক নতুন ট্যাব খুলে
\par খ) এটি হাইপারলিংক না হয়ে স্থানধারক উপাদান হিসেবে থাকে
\par গ) ফাইল ডাউনলোড হয়
\par ঘ) ছবি দেখায়
\end{minipage}
\end{enumerate}
\end{multicols}

\chFourMCQAnswer{১-খ, ২-খ, ৩-গ, ৪-খ, ৫-খ, ৬-খ, ৭-গ, ৮-খ, ৯-খ, ১০-গ, ১১-খ}

\subsubsection{গন্তব্য বোঝা যায় এমন লিংক: ছোট অনুশীলন}
শিক্ষার্থী পরীক্ষার রুটিন খুঁজছে। ``এখানে ক্লিক'' দেখে গন্তব্য বোঝা কঠিন; ``পরীক্ষার রুটিন'' নামে লিংক দিলে সিদ্ধান্ত সহজ হয়। নিচের কোডটি পেজের \texttt{body}-তে লেখা হবে।
\par\noindent\begin{minipage}[t]{.50\linewidth}\vspace{0pt}\small
\textbf{\color{chFourTeal}কোড}\par
\begin{verbatim}
<a href="exam-routine.html">
  Examination routine
</a>
\end{verbatim}
\end{minipage}\hfill
\begin{minipage}[t]{.46\linewidth}\vspace{0pt}\small
\textbf{\color{chFourTeal}আগে ভাবো, পরে পরীক্ষা করো}\par
দৃশ্যমান লেখা Admission notice করলেও গন্তব্য বদলায় না; গন্তব্য থাকে \texttt{href}-এ। লেখাটি গন্তব্যের সঙ্গে না মিললে ব্যবহারকারীর ভুল প্রত্যাশা তৈরি হয়। লিংকের ঠিকানা ও দৃশ্যমান পরিচয় দুটোই মিলতে হয়।
\end{minipage}\par


